\section{Register of the catalog}
\label{supp:register}

\subsection{The 36 entries}

\Cref{tab:register} gives each entry its identifier and name, its kind
(\Cref{sec:setting-kinds}), where the main paper discusses it, what it establishes, and the
evidence attached to it.

{\scriptsize
\begin{longtable}{@{}>{\raggedright\arraybackslash}p{3.0cm}>{\raggedright\arraybackslash}p{3.0cm}>{\raggedright\arraybackslash}p{6.2cm}>{\raggedright\arraybackslash}p{2.6cm}@{}}
\caption{Register of the catalog.}\label{tab:register}\\
\toprule
\textbf{Entry} & \textbf{Kind; place} & \textbf{What is established} & \textbf{Evidence} \\
\midrule
\endfirsthead
\toprule
\textbf{Entry} & \textbf{Kind; place} & \textbf{What is established} & \textbf{Evidence} \\
\midrule
\endhead
\bottomrule
\endfoot
VII-C001 \newline Accessible-domain state normal form & definition \newline \Cref{def:access} & Defines the seven-coordinate access record and its coherence spine. Each failed converse has a coherent witness in the reference world. & assays P5-E11; countermodels CM-08, CM-09, CM-12; 48 Lean declarations \\
VII-C002 \newline Lawful admission transition system & definition \newline \Cref{sec:certificates} & Defines five typed admission transitions, each requiring source, budget, guard and audit witnesses, with a typed effect on admissibility. & assays P5-E01; countermodels CM-08, CM-09, CM-18; 41 Lean declarations \\
VII-C003 \newline Bootstrap obstruction theorem & consequence of definitions (\Cref{prop:bootstrap}) \newline \Cref{prop:bootstrap} & In a closed regime with no admitted seed and no reachable generator no first extension is authorized; this follows from the definition of authorization. Each of the three escape routes is realized by a regime record. & assays P5-E06; countermodels CM-19, CM-20; 31 Lean declarations \\
VII-C004 \newline Neutral seed and provisioning certificate & definition; authorization does not use neutrality \newline \Cref{sec:sep-definitional} & Defines the neutral-provisioning certificate and refuses neutral credit to retrospectively stocked or relabelled system-generated seeds. Authorization of a first extension does not use neutrality. & countermodels CM-01, CM-19; 44 Lean declarations \\
VII-C005 \newline Common-origin non-transfer law & separations and a refusal \newline \Cref{sec:sep-separations} & Exhibits access records with the same origin, carrier and instrument and different access fields, and a record with shared access and a common source; a shared lineage fails the independence gate by definition. & countermodels CM-01, CM-12, CM-26; 35 Lean declarations \\
VII-C006 \newline Prospective commitment certificate & definition \newline \Cref{sec:certificates} & Defines the prospective-commitment record; a use registered after its evidence cannot discharge the prospective certificate. & assays P5-E01; countermodels CM-20, CM-09; 26 Lean declarations \\
VII-C007 \newline Join-entry record normal form & definition \newline \Cref{sec:certificates} & Defines the join-entry record. A well-formed record does not establish that a join exists. & assays P3-E01, P3-E02; countermodels CM-02, CM-03, CM-14; 32 Lean declarations \\
VII-C008 \newline Join existence and obstruction status calculus & definition \newline \Cref{def:join-status} & Defines seven join statuses; over each declared finite interface family the classifier assigns exactly one, checked by enumeration. & assays P3-E02, P5-E03, P5-E11; countermodels CM-02, CM-13, CM-27; 40 Lean declarations \\
VII-C009 \newline Strict join certificate and anti-product witness & definition; semantic counterpart \Cref{thm:bridge} \newline \Cref{def:strict-certificate} & Defines the strict-join certificate; its necessity clauses are projections of recorded fields. Semantic counterpart: main paper, \Cref{thm:strict-split} and \Cref{thm:bridge}, which cover the parents and their pairing but not every common refinement (\Cref{rem:refinement-limit}). & assays P3-E03, P3-E10, P5-E01, P5-E02, P5-E07; countermodels CM-03, CM-13, CM-14, CM-16; 44 Lean declarations \\
VII-C010 \newline Source-independence gate & definition \newline \Cref{sec:certificates} & Defines the source-independence gate; behavioural non-factorization alone does not earn independence-sensitive credit. & assays P3-E04, P5-E02; countermodels CM-01, CM-26; 28 Lean declarations \\
VII-C011 \newline Join budget and payment ledger & definition \newline \Cref{sec:certificates} & Defines the join ledger; an entry is credited when payments plus refunds cover its amount, or when its amount is zero on a certified zero-cost channel. & assays P3-E05, P3-E06, P5-E01, P5-E02, P5-E06; countermodels CM-11, CM-21; 32 Lean declarations \\
VII-C012 \newline Enablement record and attribution calculus & definition \newline \Cref{def:enablement} & Defines the enablement record with its attribution field; a hidden executor defeats attribution. & assays P4-E01, P5-E11; countermodels CM-06, CM-07, CM-17; 39 Lean declarations \\
VII-C013 \newline Endogenous enablement criterion & definition read as an equivalence \newline \Cref{sec:sep-definitional} & The endogeny criterion is the definition read as an equivalence. The two hidden-executor controls each satisfy every other condition. & assays P4-E02, P5-E09; countermodels CM-11, CM-17; 35 Lean declarations \\
VII-C014 \newline Birth/contact-surface classification & definition \newline \Cref{sec:certificates} & Defines the birth record; a contact record can be marked relation-only. & assays P4-E03; countermodels CM-02, CM-17; 42 Lean declarations \\
VII-C015 \newline Typed transmission and descent-fidelity law & consequence of the definition; semantic counterpart \Cref{prop:selection} \newline \Cref{sec:sep-definitional} & Defines transmission records; for a pure downward selection the conclusion that no lower fact is created is contained in the definition. & assays P4-E04, P5-E05; countermodels CM-06, CM-22, CM-23; 46 Lean declarations \\
VII-C016 \newline Peer contact and transport without join & separation (existence of transport without a composite) \newline \Cref{sec:sep-separations} & Exhibits a well-formed peer-transport record marked as forming no composite. & assays P3-E01; countermodels CM-02, CM-19; 42 Lean declarations \\
VII-C017 \newline Interaction-order residue and holonomy & definition; separation from the arrow \newline \Cref{def:route} & Defines route residue for two route results; exhibits a comparison record with residue and no arrow credit. & assays P4-E09, P4-E10, P5-E07, P5-E10; countermodels CM-10, CM-24; 46 Lean declarations \\
VII-C018 \newline Admission confluence and seed-dependence law & reformulation; semantic counterpart \Cref{thm:admission} \newline \Cref{sec:sep-definitional} & Relates two predicates on the same critical-pair records; this is a reformulation and proves no path confluence. Exhibits a seed-partition record whose two terminal-package fields are unequal. Semantic counterpart: main paper, \Cref{thm:admission}, which gives a unique terminal state from each state when admission only enables, and \Cref{ex:disabling}. & assays P4-E07, P4-E08, P5-E10; countermodels CM-23, CM-24; 39 Lean declarations \\
VII-C019 \newline Access/join residual and obstruction-dissolution ledger & definition \newline \Cref{sec:certificates} & Defines the append-only residual ledger with inherited, dissolved and join-created classes; scope changes do not delete residuals silently. & assays P3-E08, P5-E04, P5-E11; countermodels CM-22, CM-25; 25 Lean declarations \\
VII-C020 \newline Bridge and semantic-withdrawal discipline & definition \newline \Cref{sec:certificates} & Defines the bridge contract; a citation alone does not license transport. & countermodels CM-12, CM-18; 10 Lean declarations \\
VII-C021 \newline Reachability, guard activity, and horizon law & certificate definition; semantic counterpart \Cref{prop:search} \newline \Cref{sec:sep-definitional} & Defines the operational profile and the horizon-qualified negative certificate; reachability alone does not give occurrence. & assays P5-E08; countermodels CM-08, CM-09, CM-18; 45 Lean declarations \\
VII-C022 \newline Exposure, recoverability, admissibility, and rigidity calculus & definition \newline \Cref{sec:certificates} & Defines interface coordinates for determination, exposure, recoverability and admissibility; equal determination can have different exposure. & countermodels CM-03, CM-12; 48 Lean declarations \\
VII-C023 \newline Negative-result quantifier discipline & definition \newline \Cref{sec:certificates} & Defines the negative-evidence record; a point null does not license an unrestricted negative. & assays P5-E08; countermodels CM-18, CM-27; 5 Lean declarations \\
VII-C024 \newline Claim-grade and normative-specification discipline & definition \newline \Cref{sec:setting-kinds} & Defines the grade record; finite evidence does not upgrade a claim's grade. & countermodels CM-12, CM-18; 10 Lean declarations \\
VII-C025 \newline Two-Theory World detector contract & decidable finite statement \newline \Cref{sec:ttw} & Over the 24 scenarios and 27 named countermodels the detector contract is well formed and complete and every expected status is decided; checked by evaluation. & assays P5-E11; countermodels CM-01, CM-02, CM-03, CM-04, CM-05, CM-06, CM-07, CM-08, CM-09, CM-10, CM-11, CM-12, CM-13, CM-14, CM-15, CM-16, CM-17, CM-18, CM-19, CM-20, CM-21, CM-22, CM-23, CM-24, CM-25, CM-26, CM-27; 8 Lean declarations \\
VII-C026 \newline Categorical reduction decision & deferred; checklist only \newline \Cref{rem:deferral-categorical} & Deferred. The categorical certificate is a checklist of flags and supports no categorical statement; a mere product earns no strict-join credit. & assays P3-E10; countermodels CM-13, CM-27; 21 Lean declarations \\
VII-C027 \newline Enablement-chain composition law & associativity unconditional; credit needs composability \newline \Cref{sec:sep-definitional} & Resource deltas combine associatively with an identity for all inputs; composability governs credit. Exhibits a two-link chain record whose composed values differ with order. & assays P4-E06, P5-E04; countermodels CM-07, CM-24; 45 Lean declarations \\
VII-C028 \newline Primitive-operation algebra readiness criterion & deferred; monoid of resource deltas \newline \Cref{rem:deferral-algebra} & Deferred; the resource-delta fragment is a monoid. & assays P4-E12; countermodels CM-24; 26 Lean declarations \\
VII-C029 \newline Observer and instrument occupancy law & refusal \newline \Cref{sec:sep-refusals} & Native or endogenous credit is refused while charged occupancy is below actual occupancy; external observer and zero occupancy are escapes. & assays P5-E02, P5-E09; countermodels CM-11, CM-21; 29 Lean declarations \\
VII-C030 \newline Parent refinement, retention, and join descent & refuted; four-way classification \newline \Cref{rem:refinement-refuted} & Refuted at the level of records: a well-formed refinement record labelled destroys exists, so unconditional preservation is not permitted. All four effect labels occur among the records of the declared carrier; no model of refinement acting on a join is given. & assays P3-E07, P5-E05; countermodels CM-05, CM-22; 29 Lean declarations \\
VII-C031 \newline Contact-degree conservation decision & two scalars disagree \newline \Cref{rem:degree-refuted} & Two particular scalar degrees disagree on one vector, and an apparent conservation omitting observer occupancy fails the accounting certificate. Nothing universal is proved. & assays P3-E11; countermodels CM-21, CM-25; 22 Lean declarations \\
VII-C032 \newline Interaction arrow, irreversibility, and cross-time contact & refusal without a drive certificate \newline \Cref{def:drive} & An arrow claim requires a drive certificate; exhibits a comparison record with residue and no arrow credit. Cross-time contact requires a synchronization or partial-order witness. & assays P4-E10, P4-E11, P5-E07, P5-E10; countermodels CM-10, CM-24; 58 Lean declarations \\
VII-C033 \newline Coverage-qualified certified non-interaction & certificate definition; semantic counterpart \Cref{prop:search} \newline \Cref{sec:sep-definitional} & Defines the coverage certificate for non-interaction; a declared exhaustive family excludes contact within that family only. & assays P3-E09, P5-E03; countermodels CM-02, CM-18, CM-27; 39 Lean declarations \\
VII-C034 \newline Enablement without descent separation theorem & separation profiles \newline \Cref{sec:sep-separations} & Exhibits Boolean separation profiles: load-bearing enablement with no recorded descent, and a necessary but insufficient enabler. & assays P4-E05, P5-E05; countermodels CM-06, CM-07; 25 Lean declarations \\
VII-C035 \newline No-free-join and self-bootstrap no-go family & refusals; first clause \Cref{prop:no-free-join} \newline \Cref{sec:sep-refusals} & Payment and relabelling refusals follow from the ledger and anti-product definitions. The first clause, that a join cannot create its own first capability, is \Cref{prop:no-free-join} of the main paper. & assays P3-E05, P5-E06; countermodels CM-03, CM-19, CM-21; 53 Lean declarations \\
VII-C036 \newline Join-created obstruction and needle law & recorded attribution \newline \Cref{sec:sep-definitional} & The ledger can record a transition that discharges one inherited residual and marks one new residual as join-created. The mark is recorded, not derived from a model of a join. & assays P3-E08, P5-E04; countermodels CM-25, CM-03; 25 Lean declarations \\

\end{longtable}
}

\subsection{Recorded non-implications}

\Cref{tab:nonclaims} lists, for each entry, the non-implications recorded with it and its
escape routes.

{\scriptsize
\begin{longtable}{@{}l>{\raggedright\arraybackslash}p{8.2cm}>{\raggedright\arraybackslash}p{5.0cm}@{}}
\caption{Recorded non-implications and escape routes.}\label{tab:nonclaims}\\
\toprule
\textbf{Entry} & \textbf{Does not claim} & \textbf{Escape routes} \\
\midrule
\endfirsthead
\toprule
\textbf{Entry} & \textbf{Does not claim} & \textbf{Escape routes} \\
\midrule
\endhead
\bottomrule
\endfoot
VII-C001 & Domain width is not capability. No total order among all status coordinates is claimed. Finite reference-world success does not universalize the statement. & Narrow or change the declared carrier/interface family. Add an explicit source, seed, bridge, budget, or drive witness omitted by the no-go hypotheses. \\
VII-C002 & A well-formed rule need not be executable. A reachable transition need not fire. Finite reference-world success does not universalize the statement. & Narrow or change the declared carrier/interface family. Add an explicit source, seed, bridge, budget, or drive witness omitted by the no-go hypotheses. \\
VII-C003 & External provision is not forbidden. The statement is scoped to the declared transition family. Finite reference-world success does not universalize the statement. & Narrow or change the declared carrier/interface family. Add an explicit source, seed, bridge, budget, or drive witness omitted by the no-go hypotheses. \\
VII-C004 & Neutrality is audit-relative, not metaphysical independence. Finite reference-world success does not universalize the statement. & Narrow or change the declared carrier/interface family. Add an explicit source, seed, bridge, budget, or drive witness omitted by the no-go hypotheses. \\
VII-C005 & Common origin may still be relevant evidence when correctly typed. Finite reference-world success does not universalize the statement. & Narrow or change the declared carrier/interface family. Add an explicit source, seed, bridge, budget, or drive witness omitted by the no-go hypotheses. \\
VII-C006 & Commitment does not guarantee reachability or success. Finite reference-world success does not universalize the statement. & Narrow or change the declared carrier/interface family. Add an explicit source, seed, bridge, budget, or drive witness omitted by the no-go hypotheses. \\
VII-C007 & Record completeness does not prove a join exists. Finite reference-world success does not universalize the statement. & Narrow or change the declared carrier/interface family. Add an explicit source, seed, bridge, budget, or drive witness omitted by the no-go hypotheses. \\
VII-C008 & The taxonomy is not universal without a completeness theorem. Finite reference-world success does not universalize the statement. & Narrow or change the declared carrier/interface family. Add an explicit source, seed, bridge, budget, or drive witness omitted by the no-go hypotheses. \\
VII-C009 & Strict join does not imply an arrow. Parent strictness alone does not prove joint strictness. Finite reference-world success does not universalize the statement. & Narrow or change the declared carrier/interface family. Add an explicit source, seed, bridge, budget, or drive witness omitted by the no-go hypotheses. \\
VII-C010 & Source independence is not sufficient for strict novelty or causal independence. Finite reference-world success does not universalize the statement. & Narrow or change the declared carrier/interface family. Add an explicit source, seed, bridge, budget, or drive witness omitted by the no-go hypotheses. \\
VII-C011 & No universal conserved interaction currency is asserted. Finite reference-world success does not universalize the statement. & Narrow or change the declared carrier/interface family. Add an explicit source, seed, bridge, budget, or drive witness omitted by the no-go hypotheses. \\
VII-C012 & Enablement does not imply descent or sufficient determination. Finite reference-world success does not universalize the statement. & Narrow or change the declared carrier/interface family. Add an explicit source, seed, bridge, budget, or drive witness omitted by the no-go hypotheses. \\
VII-C013 & Endogenous does not mean uncaused or environmentally isolated. Finite reference-world success does not universalize the statement. & Narrow or change the declared carrier/interface family. Add an explicit source, seed, bridge, budget, or drive witness omitted by the no-go hypotheses. \\
VII-C014 & Contact does not itself constitute birth. Finite reference-world success does not universalize the statement. & Narrow or change the declared carrier/interface family. Add an explicit source, seed, bridge, budget, or drive witness omitted by the no-go hypotheses. \\
VII-C015 & Construction direction is not automatically temporal directionality. Finite reference-world success does not universalize the statement. & Narrow or change the declared carrier/interface family. Add an explicit source, seed, bridge, budget, or drive witness omitted by the no-go hypotheses. \\
VII-C016 & Peer transport is not joint objecthood or strict novelty. Finite reference-world success does not universalize the statement. & Narrow or change the declared carrier/interface family. Add an explicit source, seed, bridge, budget, or drive witness omitted by the no-go hypotheses. \\
VII-C017 & P3 holonomy does not imply an arrow. Finite reference-world success does not universalize the statement. & Narrow or change the declared carrier/interface family. Add an explicit source, seed, bridge, budget, or drive witness omitted by the no-go hypotheses. \\
VII-C018 & One commuting square does not prove global confluence. Finite reference-world success does not universalize the statement. & Narrow or change the declared carrier/interface family. Add an explicit source, seed, bridge, budget, or drive witness omitted by the no-go hypotheses. \\
VII-C019 & Access growth need not monotonically reduce all obstructions. Finite reference-world success does not universalize the statement. & Narrow or change the declared carrier/interface family. Add an explicit source, seed, bridge, budget, or drive witness omitted by the no-go hypotheses. \\
VII-C020 & A bridge certificate does not strengthen its source theorem. Finite reference-world success does not universalize the statement. & Narrow or change the declared carrier/interface family. Add an explicit source, seed, bridge, budget, or drive witness omitted by the no-go hypotheses. \\
VII-C021 & Non-occurrence alone is not impossibility. Finite reference-world success does not universalize the statement. & Narrow or change the declared carrier/interface family. Add an explicit source, seed, bridge, budget, or drive witness omitted by the no-go hypotheses. \\
VII-C022 & Exposure is not determination; recoverability is not permission. Finite reference-world success does not universalize the statement. & Narrow or change the declared carrier/interface family. Add an explicit source, seed, bridge, budget, or drive witness omitted by the no-go hypotheses. \\
VII-C023 & A failed instance never refutes an unrestricted existential claim. Finite reference-world success does not universalize the statement. & Narrow or change the declared carrier/interface family. Add an explicit source, seed, bridge, budget, or drive witness omitted by the no-go hypotheses. \\
VII-C024 & Finite reference-world success does not universalize the statement. & Narrow or change the declared carrier/interface family. Add an explicit source, seed, bridge, budget, or drive witness omitted by the no-go hypotheses. \\
VII-C025 & Passing finite scenarios does not prove universal laws. Finite reference-world success does not universalize the statement. & Narrow or change the declared carrier/interface family. Add an explicit source, seed, bridge, budget, or drive witness omitted by the no-go hypotheses. \\
VII-C026 & No unconditional categorical reduction is claimed. Finite reference-world success does not universalize the statement. & Narrow or change the declared carrier/interface family. Add an explicit source, seed, bridge, budget, or drive witness omitted by the no-go hypotheses. \\
VII-C027 & Enablement is not assumed transitive without hypotheses. Finite reference-world success does not universalize the statement. & Narrow or change the declared carrier/interface family. Add an explicit source, seed, bridge, budget, or drive witness omitted by the no-go hypotheses. \\
VII-C028 & No complete primitive algebra is claimed. Finite reference-world success does not universalize the statement. & Narrow or change the declared carrier/interface family. Add an explicit source, seed, bridge, budget, or drive witness omitted by the no-go hypotheses. \\
VII-C029 & Observation is not prohibited; unrecorded load-bearing observation is. Finite reference-world success does not universalize the statement. & Narrow or change the declared carrier/interface family. Add an explicit source, seed, bridge, budget, or drive witness omitted by the no-go hypotheses. \\
VII-C030 & Retention alone is not strict novelty. Finite reference-world success does not universalize the statement. & Narrow or change the declared carrier/interface family. Add an explicit source, seed, bridge, budget, or drive witness omitted by the no-go hypotheses. \\
VII-C031 & No conserved contact quantity is claimed. Finite reference-world success does not universalize the statement. & Narrow or change the declared carrier/interface family. Add an explicit source, seed, bridge, budget, or drive witness omitted by the no-go hypotheses. \\
VII-C032 & Order sensitivity and entropy production internal to a parent do not alone prove an inter-theory arrow. Finite reference-world success does not universalize the statement. & Narrow or change the declared carrier/interface family. Add an explicit source, seed, bridge, budget, or drive witness omitted by the no-go hypotheses. \\
VII-C033 & No finite null universalizes beyond its covered family. Finite reference-world success does not universalize the statement. & Narrow or change the declared carrier/interface family. Add an explicit source, seed, bridge, budget, or drive witness omitted by the no-go hypotheses. \\
VII-C034 & The theorem does not deny that some enablement also descends. Finite reference-world success does not universalize the statement. & Narrow or change the declared carrier/interface family. Add an explicit source, seed, bridge, budget, or drive witness omitted by the no-go hypotheses. \\
VII-C035 & The family permits admitted seeds, declared bridges, and proven zero-cost channels. Finite reference-world success does not universalize the statement. & Narrow or change the declared carrier/interface family. Add an explicit source, seed, bridge, budget, or drive witness omitted by the no-go hypotheses. \\
VII-C036 & Joining is not monotonically obstruction-reducing. Finite reference-world success does not universalize the statement. & Narrow or change the declared carrier/interface family. Add an explicit source, seed, bridge, budget, or drive witness omitted by the no-go hypotheses. \\

\end{longtable}
}

\subsection{Refusals}

\Cref{tab:nogo} lists the eleven refusals of the main paper's \Cref{tab:refusals} with the
entries they concern and the Lean theorems that exhibit their escape routes.

{\scriptsize
\begin{longtable}{@{}l>{\raggedright\arraybackslash}p{4.4cm}>{\raggedright\arraybackslash}p{2.6cm}>{\raggedright\arraybackslash}p{5.8cm}@{}}
\caption{Refusals and escape theorems.}\label{tab:nogo}\\
\toprule
\textbf{Identifier} & \textbf{Refusal} & \textbf{Entries} & \textbf{Escape theorems} \\
\midrule
\endfirsthead
\toprule
\textbf{Identifier} & \textbf{Refusal} & \textbf{Entries} & \textbf{Escape theorems} \\
\midrule
\endhead
\bottomrule
\endfoot
NGVII-01 & No first extension without seed or reachable generator & VII-C003 & \path{NGVII_01_escape_admitted_seed} \newline \path{NGVII_01_escape_reachable_generator} \newline \path{NGVII_01_escape_open_external_provision} \\
NGVII-02 & No resemblance-only independence-sensitive join & VII-C010, VII-C035 & \path{NGVII_02_escape_witnessed_independent_contact} \\
NGVII-03 & No retrospective self-certification & VII-C006 & \path{NGVII_03_escape_prospective_certificate} \\
NGVII-04 & No automatic total-lens transfer & VII-C005 & \path{NGVII_04_escape_certified_adapter} \\
NGVII-05 & No unpriced observer-native formation credit & VII-C029 & \path{NGVII_05_escape_external_observer} \newline \path{NGVII_05_escape_zero_occupancy} \\
NGVII-06 & No join from contact alone & VII-C007, VII-C009, VII-C035 & \path{NGVII_06_escape_complete_strict_join_certificate} \\
NGVII-07 & No product/common-refinement strictness credit & VII-C009, VII-C026, VII-C035 & \path{NGVII_07_escape_anti_product_witness} \\
NGVII-08 & No one-lineage source-independence credit & VII-C010, VII-C035 & \path{NGVII_08_escape_disjoint_root_certificate} \\
NGVII-09 & No positive-cost join credit without payment; a zero-cost join on a certified zero-cost channel is credited without payment & VII-C011, VII-C035 & \path{NGVII_09_escape_paid_channel} \newline \path{NGVII_09_escape_certified_zero_cost_channel} \\
NGVII-10 & No arrow from holonomy alone & VII-C017, VII-C032 & \path{NGVII_10_escape_independent_driven_arrow} \\
NGVII-11 & No occurrence from reachability alone & VII-C021 & \path{NGVII_11_escape_fired_occurrence} \\

\end{longtable}
}

\subsection{Cross-family corollaries}

The 21 corollaries are conjunctions and projections of certificates from different
clusters (main paper, \Cref{sec:sep-definitional}). \Cref{tab:corollaries} lists each with
its Lean theorem, the entries it spans and its location.

{\scriptsize
\begin{longtable}{@{}l>{\raggedright\arraybackslash}p{6.6cm}>{\raggedright\arraybackslash}p{2.8cm}>{\raggedright\arraybackslash}p{3.8cm}@{}}
\caption{Cross-family corollaries.}\label{tab:corollaries}\\
\toprule
\textbf{Identifier} & \textbf{Lean theorem} & \textbf{Entries} & \textbf{Location} \\
\midrule
\endfirsthead
\toprule
\textbf{Identifier} & \textbf{Lean theorem} & \textbf{Entries} & \textbf{Location} \\
\midrule
\endhead
\bottomrule
\endfoot
FVII-COR-001 & \path{certified_has_temporal_source_join_and_payment} & C002, C006, C009, C011 & \path{formalization/lean/FoundationsVII/Corollaries/ProspectiveJoin.lean}:40 \\
FVII-COR-002 & \path{retrospective_registration_blocks_prospective_join_credit} & C002, C006, C009, C011 & \path{formalization/lean/FoundationsVII/Corollaries/ProspectiveJoin.lean}:59 \\
FVII-COR-003 & \path{certified_preserves_declared_source_alignment} & C002, C006, C009, C011 & \path{formalization/lean/FoundationsVII/Corollaries/ProspectiveJoin.lean}:70 \\
FVII-COR-004 & \path{certified_has_independence_payment_and_capacity_bound} & C009, C010, C011, C029 & \path{formalization/lean/FoundationsVII/Corollaries/SourceBudget.lean}:24 \\
FVII-COR-005 & \path{source_and_payment_gates_remain_distinct} & C009, C010, C011, C029 & \path{formalization/lean/FoundationsVII/Corollaries/SourceBudget.lean}:39 \\
FVII-COR-006 & \path{certified_noninteraction_excludes_strict_join} & C008, C033 & \path{formalization/lean/FoundationsVII/Corollaries/NonInteraction.lean}:19 \\
FVII-COR-007 & \path{certified_noninteraction_retains_escape_routes} & C008, C033 & \path{formalization/lean/FoundationsVII/Corollaries/NonInteraction.lean}:35 \\
FVII-COR-008 & \path{certified_accumulates_resources_and_preserves_residuals} & C019, C027, C036 & \path{formalization/lean/FoundationsVII/Corollaries/EnablementResidual.lean}:28 \\
FVII-COR-009 & \path{join_created_needle_blocks_zero_residual_reading} & C019, C027, C036 & \path{formalization/lean/FoundationsVII/Corollaries/EnablementResidual.lean}:44 \\
FVII-COR-010 & \path{preserving_refinement_with_valid_descent_preserves_three_gates} & C015, C030, C034 & \path{formalization/lean/FoundationsVII/Corollaries/RefinementDescent.lean}:22 \\
FVII-COR-011 & \path{destroying_refinement_exhibits_a_failed_join_gate} & C015, C030, C034 & \path{formalization/lean/FoundationsVII/Corollaries/RefinementDescent.lean}:36 \\
FVII-COR-012 & \path{valid_descent_does_not_erase_refinement_obstruction} & C015, C030, C034 & \path{formalization/lean/FoundationsVII/Corollaries/RefinementDescent.lean}:45 \\
FVII-COR-013 & \path{no_free_first_join_two_gate_obstruction} & C003, C011, C035 & \path{formalization/lean/FoundationsVII/Corollaries/NoFreeJoin.lean}:12 \\
FVII-COR-014 & \path{seeded_and_paid_first_join_escape} & C003, C011, C035 & \path{formalization/lean/FoundationsVII/Corollaries/NoFreeJoin.lean}:29 \\
FVII-COR-015 & \path{strict_join_and_holonomy_do_not_replace_drive} & C009, C017, C032 & \path{formalization/lean/FoundationsVII/Corollaries/HolonomyArrow.lean}:12 \\
FVII-COR-016 & \path{holonomy_with_independent_drive_supports_arrow} & C009, C017, C032 & \path{formalization/lean/FoundationsVII/Corollaries/HolonomyArrow.lean}:26 \\
FVII-COR-017 & \path{certified_excludes_contact_and_occurrence_at_declared_scope} & C021, C023, C033 & \path{formalization/lean/FoundationsVII/Corollaries/NegativeForce.lean}:23 \\
FVII-COR-018 & \path{point_null_plus_bounded_horizon_is_not_automatically_global} & C021, C023, C033 & \path{formalization/lean/FoundationsVII/Corollaries/NegativeForce.lean}:36 \\
FVII-COR-019 & \path{certified_requires_pricing_and_internal_execution} & C013, C029 & \path{formalization/lean/FoundationsVII/Corollaries/ObserverEndogeny.lean}:22 \\
FVII-COR-020 & \path{unpriced_observer_blocks_joint_native_endogenous_credit} & C013, C029 & \path{formalization/lean/FoundationsVII/Corollaries/ObserverEndogeny.lean}:38 \\
FVII-COR-021 & \path{zero_occupancy_escape_preserves_endogenous_credit} & C013, C029 & \path{formalization/lean/FoundationsVII/Corollaries/ObserverEndogeny.lean}:48 \\

\end{longtable}
}
